\section{What the inspected $S^6$ construction actually contains}
The live workbench revision read for this tranche is
{\small\texttt{8cfbd18f3e398c533e1acabfba98b4a29975000a}}. Its S6 guide distinguishes
the original construction from later scoped reconstructions
\cite{s6repo}. We read the five-page key-advances TeX and the local-source
passages identified below, not the entire externally hosted 1,079-file,
174-MB frozen project. The source manuscript was read through its public
PDF; attempts to download it into this runtime failed DNS resolution.
No full-file hash or independent global verification is invented.

The manuscript circulated by Levent Alp\"oge describes a complex
threefold $X\to\mathbb P^1$, with complex two-torus fibres off
$0,1,\infty$. Engel's separate exposition attributes the manuscript to
Claude, communicated by Alp\"oge, and offers geometric explanation rather
than a certification asserted by this tranche \cite{source,engel}.
The following are the precise source objects, not conclusions about an
ES counterexample.

\subsection{Periods and actual duality}
The source integral matrices on $V=\mathbb L^*$ are
\[
 T_1=\begin{pmatrix}1&0&-6&2\\0&-1&1&1\\0&-1&0&1\\0&0&0&1\end{pmatrix},\quad
 T_2=\begin{pmatrix}1&6&0&-3\\0&0&-1&1\\0&1&0&0\\0&0&0&1\end{pmatrix}.
\]
The inverse-transpose relation $A_j=(T_j^{-1})^t$ proves
$\langle T_jv,A_j\ell\rangle=\langle v,\ell\rangle$ exactly.
The ordered period map is
\[
 \Pi=\begin{pmatrix}6\mu&\tau&1&0\\\beta&\mu&0&1\end{pmatrix}:\R^4\to\C^2,
 \quad \Im\tau>0,\quad
 \Im\beta-\frac{6(\Im\mu)^2}{\Im\tau}<0.
\]
These conditions make it real-linearly invertible. On the covering upper
half-plane its equivariance is
$\Pi(gz)=R_g(z)\Pi(z)\rho_V(g)^t$, with $R_g(z)\in\mathrm{GL}_2(\C)$.
The period functions obey, among others,
\[
 \tau g_1=(\tau-1)/\tau,\quad \tau g_2=-1/\tau,
 \qquad \mu g_1=(1-\mu)/\tau,\quad\mu g_2=1+\mu/\tau,
\]
\[
 \beta g_1=\beta+2-6(1-\mu)^2/\tau,\qquad
 \beta g_2=\beta-3-6\mu^2/\tau.
\]
At the two finite centres $\tau=e^{\pi i/3}$ and $i$; at the cusp the
specified $\mu$ and $\beta+\tau$ are bounded.
The common invariant covector is the $\gamma$ used in the affine slice.

On a fixed regular fibre the map
\begin{equation}\label{eq:torsion}
 (\mathbb L/D\mathbb L)\longrightarrow
 (\C^2/\Pi\mathbb L)[D],\qquad
 [v]\longmapsto[\Pi v/D]
\end{equation}
is a bijection. If $D\zeta=\Pi v$ in the torus, invertibility of $\Pi$
recovers $v$ modulo $D$; this proves both directions. Equivariance of the
periods makes \eqref{eq:torsion} commute with $A_j$.
Its restriction to $\gamma(v)=1$ is exactly the cover in this tranche.
Every such point has exact torsion order $D$.
This is a literal comparison to the regular period local system.

\subsection{Opposite-side gluing and the finite fillings}
The cusp fibre has normalization $dP_6$, the blow-up of $\mathbb P^2$ at
three noncollinear points. Its six boundary curves form a hexagon. Three
pairs of opposite curves are identified. The resulting fibre $W$ has
three double curves and two triple points, with local total-space models
$t_c=z_1$, $z_1z_2$, or $z_1z_2z_3$.
The normalization has respectively one, two, or three points over those
strata. The side-pairings are reversible on their marked branches; the
normalization-to-quotient map is not injective there. Retaining a branch
label is necessary to lift a point back.

At the finite centres the source uses ramified base coordinates
$t_j=s_j^{m_j}$, $(m_1,m_2)=(3,4)$, and free affine quotient actions
$\xi\mapsto A_j\xi+v_j/m_j$ on the real period torus, coupled with the
base rotation. The complete deck-group relations are
\[
 \Gamma_j=\langle\mathbb L,G_j\mid
 G_jt_\ell G_j^{-1}=t_{A_j\ell},\quad G_j^{m_j}=t_{v_j}\rangle.
\]
The source-side attachment kernel is the central infinite cyclic group
$\langle\widetilde G_j^{m_j}t_{-v_j}\rangle$.
The finite central product covers have degrees nine and eight, not three
and four. Their kernels and affine translations are retained in the
workbench key-advances \S3.

\subsection{The conjugate term in a period comparison}
Another literal source ``mirror'' is a real-linear comparison of nearby
periods, not a holomorphic involution. Put
$T=\Im\tau>0$, $m=\Im\mu$,
$d=\Im\beta_{c_*}-6m^2/T<0$, and
$P_*=\left(\begin{smallmatrix}0&0\\-m/T&1\end{smallmatrix}\right)$.
For $\beta_c=\beta_{c_*}+\delta$ the workbench gives
\[
 \zeta\longmapsto
 \left(I+\frac{\delta}{2id}P_*\right)\zeta
 -\frac{\delta}{2id}P_*\overline\zeta,
 \qquad \det_{\R}=\frac{d+\Im\delta}{d}>0.
\]
Both the inverse-transpose duality, the opposite-side identifications,
and this conjugate term are exact operations. They are not interchangeable.
The workbench further records a nonzero Kodaira--Spencer deformation class;
we do not reprove that global statement here.

\section{The completed filling presentation and the denominator}
The source computes a cyclic fundamental group with order
\[
 n=|12\ell_0-4\ell_1-3\ell_2|,
\]
for admissible translations, and claims that the original choice
$(\ell_0,\ell_1,\ell_2)=(0,1,-1)$ has the homology and smooth type of $S^6$.
The admissibility conditions are explicit:
$3\nmid\ell_1$ and $\ell_2$ odd \cite[Proposition 5.6 and \S7]{source}.
Their geometric justification belongs to that source. The following
presentation calculation is independent integer algebra.

\begin{proposition}[Exact cyclic presentation comparison]
For each odd $D$ prime to three, choose
\[
 \ell_0=0,\qquad
 \ell_2=\begin{cases}1&D\equiv1\ (4),\\-1&D\equiv3\ (4),\end{cases}
 \qquad \ell_1=-\frac{D+3\ell_2}{4}.
\]
These choices are admissible and have $12\ell_0-4\ell_1-3\ell_2=D$.
The relation matrix
\[
 \begin{pmatrix}-\ell_1&3&0\\-\ell_2&0&4\\-\ell_0&1&1\end{pmatrix}
\]
has Smith invariants $(1,1,D)$. Conversely every admissible nonzero value
of $12\ell_0-4\ell_1-3\ell_2$ is prime to six.
\end{proposition}
\begin{proof}
The numerator defining $\ell_1$ is divisible by four. Modulo three it is
$D$, so $3\nmid\ell_1$. The other assertions about the parameters follow
by substitution. The determinant has absolute value $D$; two minors
from the last row and the first two rows give 3 and $-4$, so the gcd of
all two-by-two minors is one, yielding the Smith invariants. Mapping the
three generators to $1$, $\ell_1/3$, $\ell_2/4$ modulo $D$ annihilates the
relations and is onto. The inverse sends $1\bmod D$ to the first generator.
An admissible value is odd modulo two and nonzero modulo three, proving
the converse. The source's original signed value is $-1$, not $+1$.
\end{proof}

Thus a denominator cyclic group can be compared to the completed-source
presentation for $\gcd(D,6)=1$. Denominators divisible by three cannot
be assigned such a free filling merely by choosing other integers in
that same $(3,4)$ family. The finite torsion-cover construction still
works for them and has the two orbits already calculated.
Neither comparison is an isomorphism between a finite ES atlas and a
compact six-manifold.

\section{A definite obstruction to one putative isomorphism}
\begin{proposition}[Genus survives the natural base change]
Take a positive-genus compact component $C_D\to\mathbb P^1$ of the cover,
and a connected resolution of a dominant component of the normalized
fibre product $X\times_{\mathbb P^1}C_D$, when the source compact fibration
is used. Its fundamental group surjects onto $\pi_1(C_D)$, so its first
Betti number is at least $2g(C_D)$. It is not diffeomorphic to $S^6$.
\end{proposition}
\begin{proof}
A loop of the compact base can be moved off the finitely many special
points. Over that complement the pullback is a proper submersion with
connected torus fibres. Lift the loop from a chosen point; connect its
endpoint back to its start within the fibre. The resulting closed loop
maps to the original loop, proving the surjection. The same open
submersion is unaffected by normalization or a resolution that is an
isomorphism over it. Abelianization gives the Betti inequality.
\end{proof}
For every odd $D>1$ at least one component has positive genus, by the
formulas already proved. This only excludes this particular pullback as
an $S^6$ copy. It does not exclude another carefully specified map, and
it is not a contradiction to ES falsity: ordinary successful primes
also have many rejected leaves with such components.

A different naive comparison already lies on a boundary. The native
coefficient form of the regular period model is
\[
 G(t,m,b)=\begin{pmatrix}2t&2m\\2m&b/3\end{pmatrix},\quad
 \det G=\frac23(tb-6m^2).
\]
Substituting the exact divisor relation $uv=a^2$ by
$(t,m,b)=(u/2,a/2,3v)$ gives
$G=\left(\begin{smallmatrix}u&a\\a&v\end{smallmatrix}\right)$ with determinant
zero. Regular periods require the strict negative determinant, so this
substitution is not a regular-fibre embedding. Our torsion comparison
uses the valid period lattice and the integer level $D$, not that
singular substitution.

\section{The level-three group, with every matrix retained}
At $D=3$ the generated matrix group has the form
\[
 \left\{
 \begin{pmatrix}1&0&0&0\\0&a&b&0\\0&c&d&0\\\chi(A)&r&s&1\end{pmatrix}:
 A=\begin{pmatrix}a&b\\c&d\end{pmatrix}\in\mathrm{SL}_2(\F_3),\quad
 r,s\in\F_3\right\}.
\]
Here $\chi:\mathrm{SL}_2(\F_3)\to\F_3$ is the character taking the planar
part of $A_1$ to 1 and that of $A_2$ to 0. For a finite exact proof,
$Q=\left(\begin{smallmatrix}0&-1\\1&0\end{smallmatrix}\right)$ and its
conjugate by $P=\left(\begin{smallmatrix}0&1\\-1&-1\end{smallmatrix}\right)$
generate the eight-element quaternion subgroup; $P$ normalizes it and
has order three. Hence $\mathrm{SL}_2(\F_3)=Q_8\rtimes\langle P\rangle$,
and the exponent of $P$ defines $\chi$.
The kernel translations with bottom rows $(0,1,0,1)$ and $(0,0,1,1)$
are respectively the source words
\[
 A_2A_2A_1A_2A_2A_1A_1,\qquad
 A_1A_1A_2A_2A_1A_2A_2 \pmod3.
\]
Multiplication verifies them. They generate all nine row translations;
together with the projected 24 matrices they give exactly the displayed
216 matrices. Removing $(r,s)$ gives an explicit splitting. The group
is therefore $(\F_3^2)^*\rtimes\mathrm{SL}_2(\F_3)$, isomorphic to
$\mathrm{ASL}_2(\F_3)$ using the invariant alternating pairing to identify
the dual module. This identification is verified as a matrix-group
statement, not inferred from the order 216 alone. No identification of
its 27-point action with the preceding Albert coordinates is asserted.
